\section{Open Markov shifts and cyclic Perron limits}
\label{sec:preliminaries}

\subsection{The killed family}

Let \(Y\subseteq A^{\mathbb N_0}\) be a topologically mixing one-step
subshift of finite type with adjacency matrix \(M\).  Let \(\varphi\)
be a one-step potential, written as \(\varphi(i,j)\) on admissible
edges, and let \(\nu_\varphi\) be its fully supported equilibrium Markov
measure.  Write \(\pi\) for the stationary row vector and \(K\) for the
transition matrix.  If \(h>0\) is a right Perron vector of
\[
  L_\varphi(i,j):=M_{ij}e^{\varphi(i,j)},
\]
then
\begin{equation}
  K_{ij}
  =M_{ij}e^{\varphi(i,j)-P_Y(\varphi)}\frac{h_j}{h_i}.
  \label{eq:markov-normalization}
\end{equation}

Choose a nonempty proper set \(A_{\rm hole}\subsetneq A\), put
\[
  H:=\bigcup_{a\in A_{\rm hole}}[a],
  \qquad
  \mathcal S:=A\setminus A_{\rm hole},
\]
and define
\[
  \tau_H(y):=\inf\{n\geq0:y_n\in A_{\rm hole}\}.
\]
Throughout the paper we make the following standing assumption.

\medskip
\noindent\textbf{Periodic survivor hypothesis.}
The graph induced by \(\mathcal S\) is irreducible of period \(d\).
Its cyclic decomposition is denoted by
\[
  \mathcal S=\mathcal C_0\sqcup\cdots\sqcup\mathcal C_{d-1},
\]
with indices in \(\mathbb Z/d\mathbb Z\), chosen so that every safe
edge from \(\mathcal C_a\) ends in \(\mathcal C_{a+1}\).

\medskip
This hypothesis implies that every safe state lies on an infinite
survivor orbit; in particular there are no transient safe states.  The
survivor subshift \(Y_H\) is irreducible, but it is mixing only when
\(d=1\).

For \(s>0\), define
\[
  u_s(i):=\pi_i^s,
  \qquad
  B_s(i,j):=K_{ij}^s\mathbf1_{\{j\in\mathcal S\}},
  \qquad i,j\in\mathcal S,
\]
and put \(\rho_s:=\rho(B_s)\).  The support of \(B_s\) is independent
of \(s\), so every \(B_s\) is irreducible of period \(d\) with the
same cyclic classes.

Let \(r_s>0\) and \(\ell_s>0\) be right and left Perron vectors,
normalized by
\[
  B_sr_s=\rho_sr_s,
  \qquad
  \ell_s^{\mathsf T}B_s=\rho_s\ell_s^{\mathsf T},
  \qquad
  \ell_s^{\mathsf T}r_s=1.
\]
For a vector \(v\) and a set \(C\), write \(v_C\) for its restriction
to \(C\).  For \(a,r\in\mathbb Z/d\mathbb Z\), set
\[
  U_{s,a}:=u_{s,\mathcal C_a}^{\mathsf T}r_{s,\mathcal C_a},
  \qquad
  V_{s,a}(v):=\ell_{s,\mathcal C_a}^{\mathsf T}v_{\mathcal C_a},
\]
and define the cyclic coefficient
\begin{equation}
  A_{s,r}(v)
  :=d\sum_{a\in\mathbb Z/d\mathbb Z}
       U_{s,a}V_{s,a+r}(v).
  \label{eq:cyclic-coefficient}
\end{equation}
If \(v>0\), then \(A_{s,r}(v)>0\) for every \(r\).

\subsection{Cyclic asymptotics}

We give the precise matrix limit used throughout.  Although it is a
standard periodic Perron formula, keeping the coefficient visible is
essential: different entrywise powers \(B_s\) have different Perron
vectors, so normalization by the survivor mass does not cancel the
depth phase.

\begin{proposition}[Cyclic Perron coefficient]
\label{prop:cyclic-perron}
Fix \(s>0\) and a vector \(v\).  For each
\(r\in\mathbb Z/d\mathbb Z\), there is \(0\leq\kappa_{s,r}<1\) such
that, as \(n\to\infty\) through \(n\equiv r\pmod d\),
\begin{equation}
  u_s^{\mathsf T}B_s^n v
  =\rho_s^n\bigl(A_{s,r}(v)+O(\kappa_{s,r}^n)\bigr).
  \label{eq:cyclic-perron-asymptotic}
\end{equation}
The error is uniform over any fixed finite collection of vectors \(v\).
\end{proposition}

\begin{proof}
Normalize \(B_s\) to the stochastic matrix
\[
  P_s(i,j):=\frac{B_s(i,j)r_s(j)}{\rho_s r_s(i)}.
\]
Its stationary distribution is
\[
  \zeta_s(i):=\ell_s(i)r_s(i),
\]
because \(\ell_s^{\mathsf T}r_s=1\).  Stationarity and the cyclic
decomposition imply
\[
  \sum_{i\in\mathcal C_a}\zeta_s(i)=\frac1d
  \qquad(a\in\mathbb Z/d\mathbb Z).
\]
Indeed, one-step probability flow transports the total stationary mass of
\(\mathcal C_a\) to \(\mathcal C_{a+1}\), so all class masses are
equal.

For each \(a\), the restriction of \(P_s^d\) to
\(\mathcal C_a\) is primitive.  The primitive Perron theorem therefore
gives, for \(i\in\mathcal C_a\) and
\(j\in\mathcal C_{a+r}\),
\[
  P_s^n(i,j)
  =d\zeta_s(j)+O(\kappa_{s,r}^n)
  \qquad(n\equiv r\pmod d),
\]
while \(P_s^n(i,j)=0\) if \(j\notin\mathcal C_{a+r}\).  The state
space is finite, so the error may be chosen uniformly in \(i,j\).

Diagonal similarity gives
\[
  B_s^n(i,j)
  =\rho_s^n\frac{r_s(i)}{r_s(j)}P_s^n(i,j).
\]
Since \(\zeta_s(j)/r_s(j)=\ell_s(j)\), summing against
\(u_s(i)v(j)\) yields
\begin{align*}
  \rho_s^{-n}u_s^{\mathsf T}B_s^n v
  &=d\sum_a
      \left(\sum_{i\in\mathcal C_a}u_s(i)r_s(i)\right)
      \left(\sum_{j\in\mathcal C_{a+r}}\ell_s(j)v(j)\right)
      +O(\kappa_{s,r}^n)\\
  &=A_{s,r}(v)+O(\kappa_{s,r}^n),
\end{align*}
which proves the claim.
\end{proof}

\begin{remark}[Peripheral-spectrum form]
The limit in Proposition~\ref{prop:cyclic-perron} is equivalently the
sum of the \(d\) peripheral spectral projections of \(B_s\), weighted
by the appropriate roots of unity.  Formula
\eqref{eq:cyclic-coefficient} is preferable here because it is positive,
real, and directly computable from Perron vectors.  See
\citet{Seneta2006} for the periodic theory of irreducible nonnegative
matrices.
\end{remark}

\subsection{Safe words and pressure}

Let \(\mathcal L_m^H\) be the set of admissible words
\(a=a_0\cdots a_{m-1}\) whose symbols all lie in \(\mathcal S\), and
define
\[
  Z_{H,m}(s):=\sum_{a\in\mathcal L_m^H}\nu_\varphi([a])^s.
\]
The one-step Markov property gives the exact path sum
\begin{equation}
  Z_{H,m}(s)=u_s^{\mathsf T}B_s^{m-1}\mathbf1.
  \label{eq:power-sum-path}
\end{equation}
In particular,
\[
  Z_{H,m}(1)=\nu_\varphi(\tau_H\geq m).
\]

For later use we identify \(\rho_s\) thermodynamically.  Let
\(D_s\) be diagonal with \(D_s(i,i)=h_i^s\), where \(h\) is the vector
in \eqref{eq:markov-normalization}.  On \(\mathcal S\),
\[
  B_s
  =e^{-sP_Y(\varphi)}D_s^{-1}L_{s\varphi,H}D_s,
\]
where
\[
  L_{s\varphi,H}(i,j)
  :=M_{ij}e^{s\varphi(i,j)}\mathbf1_{\{j\in\mathcal S\}}.
\]
The survivor graph is irreducible, so the usual pressure formula remains
valid without primitivity:
\begin{equation}
  \rho_s
  =\exp\bigl(P_{Y_H}(s\varphi)-sP_Y(\varphi)\bigr).
  \label{eq:rho-pressure}
\end{equation}
Consequently the escape exponent is
\[
  -\log\rho_1=P_Y(\varphi)-P_{Y_H}(\varphi).
\]
Periodicity affects the prefactor of survivor mass, but not its
exponential rate.

Combining \eqref{eq:power-sum-path} with
Proposition~\ref{prop:cyclic-perron} gives the basic phase-resolved
formula
\begin{equation}
  Z_{H,m}(s)
  =A_{s,r}(\mathbf1)\rho_s^{m-1}
    (1+O(\kappa_s^{m-1})),
  \qquad m-1\equiv r\pmod d,
  \label{eq:unconditioned-phase-power}
\end{equation}
where, after taking a maximum over the finitely many phases, one may use
a single \(0<\kappa_s<1\).
